%%%%%%%%%%%%%%%%%%%%%%%%%%%%%%%%%%%%%%%%%%%%%%%%%%%%%%%%%%%%%%%%%%%%%%%%%%%%%%%%
\section{Cross-State Inference (RQ5)}
\label{sec:results:rq5}
% Answers \ref{sec:intro:rq5}: train on reduced, infer on full.

\subsection{Reduced-to-Full Inference Accuracy}
\label{sec:results:rq5:accuracy}
Smooth L1 loss, RMSE and $R^2$ from querying each reduction-trained model on unreduced AIGs,
against the same model scored under the reduction it was trained with
(\ref{sec:method:experiment:metrics}).

\begin{figure}[htbp]
    \centering
    \includegraphics[width=\columnwidth]{\resultsfigs/rq5_positive_control}
    \caption[RQ5 positive control]{\textbf{ENTIRELY FABRICATED.} Colour refinement
    at count-cap $\infty$ merges only nodes the encoder provably cannot
    distinguish, so a model trained on those graphs \emph{must} land on the dashed
    baseline when queried on full graphs. Without this control a poor transfer
    elsewhere cannot be told apart from a broken evaluation path, which is why the
    slot is held open rather than dropped.}
    \label{fig:rq5_positive_control}
\end{figure}

The control is reported first because it bounds every other result in this section
(\ref{sec:method:architecture:exact}): if it lands away from the dashed baseline once
measured, cross-state failure elsewhere cannot be separated from a defect in the evaluation
path itself.

\begin{figure*}[htbp]
    \centering
    \includegraphics[width=\textwidth]{\resultsfigs/rq5_matched_vs_full}
    \caption[Matched-state against full-graph evaluation]{The same weights
    evaluated under the reduction they were trained with and on unreduced graphs.
    For every method here a full graph needs no conversion, since it already is a
    valid input, which is what makes this a direct test rather than a comparison
    across two input formats (\ref{sec:method:summary:schema}).}
    \label{fig:rq5_matched_vs_full}
\end{figure*}

\begin{table*}[htbp]
\centering
\caption{RQ5: the same weights evaluated under the reduction they were trained with and on unreduced graphs. For every method here a full graph needs no conversion, it already is a valid input. The exact colour-refinement track is the exception and has not been run.}
\label{tab:rq5_cross_state}
\small
\resizebox{\textwidth}{!}{%
\begin{tabular}{llrrrrrr}
\toprule
\textbf{Method} & \textbf{Family} & \textbf{$R^2$ matched} & \textbf{$R^2$ full} & \textbf{$\Delta R^2$} & \textbf{$\rho$ matched} & \textbf{$\rho$ full} & \textbf{Throughput full (g/s)} \\
\midrule
Full graph & Baseline & -- & \textbf{0.343} & -- & -- & 0.185 & 74.650 \\
Random hashing & Partition & -0.981 & -0.071 & 0.910 & -0.183 & 0.145 & 74.669 \\
METIS & Partition & -3.062 & -2.663 & 0.399 & 0.237 & 0.225 & 74.668 \\
Level slicing & Partition & -0.142 & -0.142 & 0.000 & 0.119 & \textbf{0.601} & 74.690 \\
Span-weighted METIS & Partition & 0.008 & -0.352 & -0.360 & 0.263 & -0.132 & 74.690 \\
Random edge dropout & Sparsification & \textbf{0.215} & \underline{0.247} & 0.032 & 0.194 & 0.326 & 74.658 \\
Spanning forest & Sparsification & \underline{0.138} & 0.038 & -0.100 & \underline{0.282} & \underline{0.411} & 74.721 \\
PageRank pruning & Sparsification & -0.018 & -0.114 & -0.097 & \textbf{0.367} & 0.004 & 74.721 \\
AND-gate only & Sparsification & -0.064 & -0.144 & -0.080 & -0.025 & -0.467 & 74.684 \\
\bottomrule
\end{tabular}
}
\\[4pt]\begin{minipage}{\linewidth}\footnotesize\emph{\textbf{Bold} marks the best value in a column and \underline{underline} the second best, over measured rows only. Lower is better for error, memory, time and offline cost; higher is better for $R^2$, correlation and savings.}\end{minipage}
\end{table*}


No method comes close to the full-graph baseline's $R^2 = 0.343$ once queried on full graphs.
Random edge dropout is the strongest reduction in both states, $R^2 = 0.215$ matched and
$0.247$ full, and the only sparsifier that improves on the shift at all. Random hashing moves
by the largest margin of any method, $R^2 = -0.981$ rising to $-0.071$, though neither score
is close to competitive.
\clearpage

\subsection{Zero-Shot Ranking}
\label{sec:results:rq5:ranking}
Spearman correlation under the same matched-against-full comparison, which is what the
script-selection use case in~\ref{sec:intro:motivation} depends on rather than on absolute
error.

\begin{figure}[htbp]
    \centering
    \fbox{\parbox[c][0.24\textheight][c]{0.86\columnwidth}{\centering\footnotesize
    Reserved slot. No plot generated.}}
    \caption[Zero-shot ranking under cross-state transfer]{\textbf{FABRICATED.} Reserved slot
    for the ranking view of the RQ5 comparison: Spearman correlation matched against full for
    each method, with a bootstrap interval per point. Nothing is plotted and no data stands
    behind this float. The numbers it would draw are the $\rho$ columns
    of~\ref{tab:rq5_cross_state}.}
    \label{fig:rq5_ranking}
\end{figure}

Ranking moves independently of error here just as it does under matched-state evaluation
(\ref{sec:results:rq3:ranking}). Level slicing is the extreme case: $\rho = 0.119$ matched
rises to $\rho = 0.601$ on full graphs, the best ranking score of any configuration in either
state, while its own $R^2$ stays flat and negative. AND-gate-only moves the opposite way,
from $\rho = -0.025$ matched to $-0.467$ full, ordering circuits backwards on exactly the
graphs it is meant to serve at inference time.
\clearpage

\subsection{Matched-State against Cross-State Drop-Off}
\label{sec:results:rq5:dropoff}
The cross-state drop-off $\Delta_{\mathrm{state}}$ of~\eqref{eq:method:dropoff}, the
difference the structural shift alone makes once a model is fixed.

\begin{figure}[htbp]
    \centering
    \includegraphics[width=\columnwidth]{\resultsfigs/rq5_transfer_quadrant}
    \caption[Transfer quadrant]{Matched-state against full-graph $R^2$. The
    diagonal is perfect transfer. Below it the structural shift cost accuracy,
    above it the reduction was the thing doing damage.}
    \label{fig:rq5_transfer_quadrant}
\end{figure}

Three of the eight measured configurations sit strictly above the diagonal, scoring
\emph{better} on full graphs than on the reduction they trained with, which
under~\eqref{eq:method:dropoff} reads as $\Delta_{\mathrm{state}} > 0$ and says the
reduction, not the structural shift, was costing accuracy for those methods. A fourth, level
slicing, lands exactly on it. Span-weighted METIS is the sharpest case in the other
direction, $R^2 = 0.008$ matched falling to $-0.352$ full: a configuration that looked
competitive under the condition it trained on and is not competitive under the condition RQ5
asks about.
\clearpage

\subsection{Cross-State Transfer by Reduction Family}
\label{sec:results:rq5:byfamily}
The same drop-off grouped by family, which asks whether node-preserving methods transfer
differently from node-removing ones or from partitioning as a whole
(\ref{sec:prelim:reduction}).

\begin{figure*}[htbp]
    \centering
    \includegraphics[width=\textwidth]{\resultsfigs/rq5_dropoff}
    \caption[Cost of the structural shift]{Full-graph score minus matched-state
    score, by method and family. A positive bar means the model scores
    \emph{better} on graphs it never trained on, which would say the reduction was
    hurting the model rather than the evaluation.}
    \label{fig:rq5_dropoff}
\end{figure*}

Sparsification keeps every node and changes only which edges survive, so a sparsified graph
and its full counterpart differ least as inputs. The four sparsifiers cluster near zero
drop-off, and only random edge dropout crosses into positive territory. Partitioning changes
the batch structure itself, one block at a time under matched-state and one connected graph
under full, and its four members swing wider in both directions, from span-weighted METIS's
$-0.360$ to random hashing's $+0.910$. Summarization contributes no row: no configuration in
that family has been trained, so it cannot be placed on this axis yet.
\clearpage

\subsection{Inference Cost}
\label{sec:results:rq5:cost}
Throughput and peak memory for a forward pass alone, which stores neither gradients nor
activations and so is not bound by the memory ceiling that forced training-time reduction in
the first place (\ref{sec:method:experiment:hardware}).

\begin{figure*}[htbp]
    \centering
    \includegraphics[width=\textwidth]{\resultsfigs/rq5_inference_cost}
    \caption[Inference cost by evaluation state]{Throughput and peak memory for a
    query, under each evaluation state. Inference stores neither gradients nor
    activations, so the memory argument that forced the reduction at training time
    need not bind here at all, which is the practical claim
    of~\ref{sec:intro:rq5}.}
    \label{fig:rq5_inference_cost}
\end{figure*}

Full-graph throughput sits within a tenth of a percent across every configuration
in~\ref{tab:rq5_cross_state}, from $74.650$ graphs per second for the unreduced baseline to
$74.721$ for spanning forest and PageRank pruning, so which reduction a model trained under
has no measurable effect on how fast it serves full graphs at inference time. Training-time
memory savings (\ref{sec:results:rq2:memory}) consequently do not need to be reproduced at
inference for the train-cheap, infer-full claim to hold.

\begin{figure}[htbp]
    \centering
    \includegraphics[width=\columnwidth]{\resultsfigs/rq5_cpu_inference}
    \caption[GPU against CPU inference]{\textbf{MOSTLY FABRICATED.} The ``infer
    full on modest hardware'' half of RQ5. Every surviving inference measurement
    was recorded on GPU. The CPU bars are placeholders for a CPU inference pass
    that has not been run.}
    \label{fig:rq5_cpu_inference}
\end{figure}

Modest hardware is the deployment target the practical claim of~\ref{sec:intro:rq5} is
actually about, and GPU throughput does not stand in for it: every figure above
was measured on the same device class training used, so the CPU pass is the one measurement
this chapter cannot yet substitute around.
\clearpage
